\documentclass[aspectratio=169,11pt]{beamer}

% ============================================================
% Theme and typography
% ============================================================

\usetheme{Madrid}
\usecolortheme{default}
\usefonttheme{professionalfonts}

\setbeamertemplate{navigation symbols}{}
\setbeamertemplate{headline}{}

\setbeamertemplate{footline}{%
  \leavevmode\hbox{%
  \begin{beamercolorbox}
    [wd=.84\paperwidth,ht=2.7ex,dp=1.1ex,leftskip=.8em]
    {author in head/foot}
    CSCI 6461 --- Computer Systems Architecture
  \end{beamercolorbox}%
  \begin{beamercolorbox}
    [wd=.16\paperwidth,ht=2.7ex,dp=1.1ex,rightskip=.8em]
    {date in head/foot}
    \hfill\insertframenumber/\inserttotalframenumber
  \end{beamercolorbox}}
}

\setbeamersize{
  text margin left=0.65cm,
  text margin right=0.65cm
}

\setbeamerfont{frametitle}{
  size=\Large,
  series=\bfseries
}

\setbeamerfont{block title}{
  series=\bfseries
}

\setbeamertemplate{blocks}[rounded][shadow=false]

% ============================================================
% Packages
% ============================================================

\usepackage[T1]{fontenc}
\usepackage{lmodern}
\usepackage{amsmath}
\usepackage{booktabs}
\usepackage{array}
\usepackage{tikz}
\usepackage{listings}

\usetikzlibrary{
  arrows.meta,
  positioning,
  shapes.geometric,
  calc
}

% ============================================================
% Colors
% ============================================================

\definecolor{navy}{RGB}{24,43,68}
\definecolor{deepblue}{RGB}{45,88,130}
\definecolor{lightblue}{RGB}{229,238,248}
\definecolor{softgray}{RGB}{245,247,250}
\definecolor{darkgray}{RGB}{38,43,50}
\definecolor{gold}{RGB}{190,140,40}

\setbeamercolor{palette primary}{
  bg=navy,
  fg=white
}

\setbeamercolor{palette secondary}{
  bg=deepblue,
  fg=white
}

\setbeamercolor{frametitle}{
  bg=navy,
  fg=white
}

\setbeamercolor{title}{
  fg=white
}

\setbeamercolor{subtitle}{
  fg=white
}

\setbeamercolor{author}{
  fg=white
}

\setbeamercolor{date}{
  fg=white
}

\setbeamercolor{titlelike}{
  fg=white
}

\setbeamercolor{block title}{
  bg=deepblue,
  fg=white
}

\setbeamercolor{block body}{
  bg=lightblue,
  fg=darkgray
}

\setbeamercolor{normal text}{
  fg=darkgray,
  bg=white
}

% ============================================================
% General formatting
% ============================================================

\setlength{\parskip}{0.35em}

\newcommand{\key}[1]{\textbf{#1}}
\newcommand{\bits}[1]{\texttt{#1}}
\newcommand{\hex}[1]{\texttt{0x#1}}
\newcommand{\code}[1]{\texttt{#1}}
\newcommand{\iron}{\(T_{CPU}=IC\times CPI\times T_{cycle}\)}

% ============================================================
% Metadata
% ============================================================

\title[Topic 2]{
  RISC vs. CISC:\\
  Instruction Set Design Principles \& Trade-offs
}

\subtitle{
  ISA choices, implementation costs,
  and performance consequences
}

\author{
  Computer Systems Architecture -- ARM AArch64
}

\date{}

% ============================================================
% Document
% ============================================================

\begin{document}

% ============================================================
% Slide 1
% ============================================================

\begin{frame}[plain]

\begin{tikzpicture}[remember picture,overlay]

\fill[navy]
  (current page.south west)
  rectangle
  (current page.north east);

\fill[gold]
  (current page.south west)
  rectangle
  ([yshift=.12cm]current page.south east);

\end{tikzpicture}

\vspace{1.0cm}

\begin{minipage}{0.86\textwidth}

{\color{gold}
 \large
 CSCI 6461 \;|\; TOPIC 2
 \par}

\vspace{0.28cm}

{\color{white}
 \fontsize{25}{30}\selectfont
 \bfseries
 RISC vs. CISC:\\[0.08cm]
 Instruction Set Design\\[0.08cm]
 Principles \& Trade-offs
 \par}

\vspace{0.35cm}

{\color{white!78}
 \large
 ISA choices, implementation costs,
 and performance consequences
 \par}

\end{minipage}

\vfill

{\color{white}
 \normalsize
 Dr. J. Burns
 \par}

{\color{white!65}
 \small
 Computer Systems Architecture
 \par}

\vspace{0.45cm}

\end{frame}


% 2
\begin{frame}{What Is an Instruction Set Architecture?}
\begin{itemize}
\item An ISA is the programmer-visible specification of a processor.
\item It defines instructions, registers, data types, addressing modes, and architectural state.
\item Compilers generate code for the ISA; processors implement the ISA.
\item ISA design therefore sits directly at the hardware/software boundary.
\item A good ISA must remain useful for decades while implementations evolve underneath it.
\item RISC and CISC are different historical answers to the question: \key{what should the ISA expose?}
\end{itemize}
\end{frame}

% 3
\begin{frame}{Why ISA Design Matters}
\begin{columns}[T,onlytextwidth]
\column{0.49\textwidth}
\begin{block}{Software consequences}
\begin{itemize}
\item Instruction count
\item Code size
\item Compiler complexity
\item Calling conventions
\item Binary compatibility
\end{itemize}
\end{block}
\column{0.49\textwidth}
\begin{block}{Hardware consequences}
\begin{itemize}
\item Decode complexity
\item Datapath organization
\item Pipeline regularity
\item Control logic
\item Energy per instruction
\end{itemize}
\end{block}
\end{columns}
\vspace{0.4em}
\begin{flushleft}\key{ISA design redistributes complexity between software and hardware rather than eliminating it.}\end{flushleft}
\end{frame}

% 4
\begin{frame}{The Instruction-Set Design Problem}
\begin{itemize}
\item Should one instruction perform a small primitive operation or a larger sequence of work?
\item Should arithmetic instructions access memory directly or operate only on registers?
\item Should instructions have one fixed length or many lengths?
\item How many addressing modes should hardware understand?
\item Should frequently used operations receive special instructions even if they complicate decode?
\item RISC and CISC emerged from different answers shaped by the technology of their eras.
\end{itemize}
\end{frame}

% 5
\begin{frame}{Historical Constraints Shaped Early ISAs}
\begin{itemize}
\item Early main memory was expensive, slow, and physically limited.
\item Compilers were less sophisticated than modern optimizing compilers.
\item Hardware designers wanted programs to occupy fewer bytes and make fewer memory accesses.
\item Rich instructions could express substantial work in a compact encoding.
\item Microcode made it practical to implement complex instruction behavior using simpler internal steps.
\item These pressures strongly encouraged what later became known as the CISC style.
\end{itemize}
\end{frame}

% 6
\begin{frame}{The Origins of CISC}
\begin{itemize}
\item CISC stands for \key{Complex Instruction Set Computer}.
\item The philosophy favored rich instruction sets with many operations and addressing modes.
\item Individual instructions could perform multi-step tasks such as memory access plus arithmetic.
\item Variable-length encodings allowed common instructions to be compact while retaining expressive instructions.
\item The goal was often to reduce program size and narrow the semantic gap between languages and machine code.
\item The x86 ISA is the most prominent surviving example of a CISC lineage.
\end{itemize}
\end{frame}

% 7
\begin{frame}{The CISC Design Philosophy}
\begin{itemize}
\item Provide many instructions so software can request specialized operations directly.
\item Permit operands to reside in registers or memory in many instruction forms.
\item Support multiple addressing modes to make data structures convenient to access.
\item Use variable instruction lengths to encode simple cases compactly and complex cases when needed.
\item Allow one architectural instruction to correspond to several internal hardware actions.
\item Accept more complicated decoding and control in exchange for expressive machine instructions.
\end{itemize}
\end{frame}

% 8
\begin{frame}{CISC Example: More Work per Instruction}
\begin{block}{Conceptual example}
Imagine an ISA instruction that computes \code{A = A + memory[B + 8]}.
\end{block}
\begin{itemize}
\item One architectural instruction identifies an address, reads memory, performs arithmetic, and writes a result.
\item The instruction may encode an operation, base register, displacement, and destination.
\item Fewer instructions are visible in the program.
\item Internally, however, hardware may still perform address generation, memory access, addition, and write-back separately.
\item The architectural instruction is compact semantically even if its implementation is multi-step.
\item This distinction between \key{architectural instruction} and \key{internal work} matters throughout the lecture.
\end{itemize}
\end{frame}

% 9
\begin{frame}{Why Complex Instructions Can Be Attractive}
\begin{itemize}
\item Fewer instructions can reduce instruction-fetch traffic.
\item Rich addressing modes can express common data-access patterns directly.
\item Compact encodings improve code density, which can benefit instruction-cache behavior.
\item Some specialized operations map naturally to common software idioms.
\item Backward compatibility can preserve huge software ecosystems over many processor generations.
\item The central attraction is not ``complexity for its own sake'' but doing useful work with fewer architectural instructions or bytes.
\end{itemize}
\end{frame}

% 10
\begin{frame}{The Cost of Instruction Complexity}
\begin{itemize}
\item Variable instruction length makes instruction boundaries harder to find quickly.
\item Many addressing modes increase decoder and control complexity.
\item Instructions may require very different execution times and hardware resources.
\item Irregular operand rules complicate pipelines and dependency handling.
\item Complex decode consumes area and energy before useful arithmetic even begins.
\item As transistor budgets and compiler quality improved, architects questioned whether this complexity still paid for itself.
\end{itemize}
\end{frame}

% 11
\begin{frame}{The Emergence of RISC}
\begin{itemize}
\item RISC stands for \key{Reduced Instruction Set Computer}.
\item Researchers observed that compiled programs used some instructions far more often than others.
\item Many complex instructions were rarely generated by compilers or were not faster than simple sequences.
\item Faster semiconductor technology made simpler hardwired control increasingly attractive.
\item The idea was to optimize the common case rather than provide hardware for every conceivable operation.
\item Berkeley RISC and Stanford MIPS helped establish this approach in the early 1980s.
\end{itemize}
\end{frame}

% 12
\begin{frame}{The RISC Design Philosophy}
\begin{itemize}
\item Keep instructions relatively simple and regular.
\item Emphasize register-to-register arithmetic.
\item Access memory explicitly through load and store instructions.
\item Use relatively regular instruction encodings to simplify fetch and decode.
\item Provide enough registers that temporary values can remain close to the execution units.
\item Rely more heavily on compilers to combine simple instructions into efficient programs.
\end{itemize}
\end{frame}

% 13
\begin{frame}{Berkeley RISC and Stanford MIPS}
\begin{columns}[T,onlytextwidth]
\column{0.49\textwidth}
\begin{block}{Berkeley RISC}
\begin{itemize}
\item Regular instructions
\item Large register emphasis
\item Pipeline-friendly design
\item Influenced SPARC
\end{itemize}
\end{block}
\column{0.49\textwidth}
\begin{block}{Stanford MIPS}
\begin{itemize}
\item Simple load/store ISA
\item Compiler-oriented design
\item Pipeline simplicity
\item Influential teaching model
\end{itemize}
\end{block}
\end{columns}
\vspace{0.4em}
The lasting contribution was a methodology: \key{measure real workloads, simplify the common path, and let compilers schedule simple primitives.}
\end{frame}

% 14
\begin{frame}{RISC Principle 1: Simple, Regular Operations}
\begin{itemize}
\item Regular operations reduce the number of special cases hardware must handle.
\item Similar instruction formats allow common decode circuitry to be reused.
\item Predictable operand placement simplifies datapath control.
\item Shorter critical control paths can support higher clock rates or lower energy.
\item Simplicity also makes pipelining easier because stages perform more uniform work.
\item ``Simple'' does not mean weak: complex programs are built by composing many simple instructions.
\end{itemize}
\end{frame}

% 15
\begin{frame}{RISC Principle 2: Load/Store Architecture}
\begin{block}{Key rule}
Arithmetic operates on registers; separate instructions move data between registers and memory.
\end{block}
\begin{itemize}
\item A load computes an address and brings a value into a register.
\item Arithmetic then consumes register operands.
\item A store writes a register value back to memory when needed.
\item This separates memory access from ALU computation in the ISA.
\item The separation maps naturally onto pipeline stages such as EX and MEM.
\item AArch64 follows this load/store model closely.
\end{itemize}
\end{frame}

% 16
\begin{frame}{RISC Principle 3: Register-Based Computation}
\begin{itemize}
\item Registers are much faster to access than main memory.
\item Keeping active values in registers reduces expensive memory traffic.
\item Register operands are easy to name with compact fields in an instruction.
\item Multiple read and write ports can feed execution units efficiently.
\item Compiler register allocation becomes important because software decides which values stay resident.
\item This is a deliberate transfer of responsibility: simpler hardware requires smarter code generation.
\end{itemize}
\end{frame}

% 17
\begin{frame}{RISC Principle 4: Regular Instruction Encoding}
\begin{itemize}
\item Fixed or highly regular instruction lengths make boundaries easy to locate.
\item Decoders can examine predictable bit positions for opcode and register fields.
\item Instruction fetch can retrieve aligned groups without first discovering varying lengths.
\item Regularity reduces front-end control complexity.
\item The trade-off is that some operations may require more instruction bits or multiple instructions.
\item AArch64 uses fixed 32-bit instruction encodings, a classic RISC-friendly property.
\end{itemize}
\end{frame}

% 18
\begin{frame}{RISC Principle 5: Compiler-Friendly Design}
\begin{itemize}
\item RISC assumes compilers can select and schedule simple operations effectively.
\item The compiler exposes instruction-level opportunities rather than hiding work inside large instructions.
\item Register allocation is essential because the ISA expects frequent register use.
\item Regular instructions simplify cost modeling during optimization.
\item Compiler improvements can therefore yield performance gains without changing hardware.
\item The risk is that poor code generation can produce more instructions and more memory traffic.
\end{itemize}
\end{frame}

% 19
\begin{frame}{Worked Example: Evaluate \texorpdfstring{$z=(a+b)\times c$}{z=(a+b)*c}}
\begin{block}{Load/store-style sequence}
\code{ldr x1, [a]}\\
\code{ldr x2, [b]}\\
\code{add x3, x1, x2}\\
\code{ldr x4, [c]}\\
\code{mul x5, x3, x4}\\
\code{str x5, [z]}
\end{block}
\begin{itemize}
\item The sequence is longer than a hypothetical memory-to-memory expression instruction.
\item But each instruction performs a regular, easily identified job.
\item This regularity is valuable for decoding, pipelining, dependency tracking, and compiler scheduling.
\end{itemize}
\end{frame}

% 20
\begin{frame}{Hypothetical CISC Equivalent: Same Expression}
\begin{block}{One richer architectural instruction}
\centering
\Large\code{muladd\_mem [z], [a], [b], [c]}
\end{block}
\begin{itemize}
\item This \key{hypothetical} instruction represents the same expression $z=(a+b)\times c$ in one architectural instruction.
\item The instruction must obtain three source values from memory and ultimately write one result back to memory.
\item Internally, the processor still needs address generation, multiple memory accesses, addition, multiplication, and a store.
\item Architectural instruction count falls from six visible RISC-style instructions to one visible CISC-style instruction.
\item The trade-off is a richer encoding and more complicated decode, sequencing, exception handling, and dependency behavior.
\item The example illustrates CISC philosophy; it is \key{not} intended to claim that a real x86 instruction has exactly this form.
\end{itemize}
\end{frame}

% 21
\begin{frame}{RISC vs. CISC: Same Computation, Different Exposure}
\begin{columns}[T,onlytextwidth]
\column{0.48\textwidth}
\begin{block}{RISC-style exposure}
\begin{itemize}
\item Load operands
\item Compute in registers
\item Store result
\item More visible instructions
\item More regular operations
\end{itemize}
\end{block}
\column{0.48\textwidth}
\begin{block}{CISC-style exposure}
\begin{itemize}
\item Memory may be an operand
\item Richer single instructions
\item Fewer visible instructions
\item Variable work per instruction
\item More decode/control complexity
\end{itemize}
\end{block}
\end{columns}
\vspace{0.5em}
\key{Neither description alone predicts performance.} We need the Iron Law.
\end{frame}

% 21
\begin{frame}{Instruction Count: Is Fewer Always Better?}
\begin{itemize}
\item A CISC program may require fewer architectural instructions for the same task.
\item But instruction count says nothing about the time required by each instruction.
\item A rich instruction may decode into several internal operations.
\item A RISC program may use more instructions but execute them with high pipeline regularity.
\item Code size depends on instruction width as well as instruction count.
\item Therefore \key{fewer instructions is useful only when the other performance factors remain favorable.}
\end{itemize}
\end{frame}

% 22
\begin{frame}{CPI: The Other Side of the Trade-off}
\begin{itemize}
\item CPI is the average number of clock cycles per executed architectural instruction.
\item A complex instruction may require more internal work and potentially more cycles.
\item Regular instructions can make average execution behavior easier to pipeline.
\item However, modern superscalar processors can complete multiple instructions per cycle, giving CPI below 1.
\item CPI therefore reflects implementation, workload, cache behavior, and instruction mix.
\item RISC versus CISC cannot be evaluated by instruction count without considering CPI.
\end{itemize}
\end{frame}

% 23
\begin{frame}{Clock Cycle Time and Instruction Complexity}
\begin{itemize}
\item The clock period must accommodate the critical timing requirements of the implementation.
\item More complicated decode and control can lengthen front-end paths or require additional pipeline stages.
\item Simpler instruction structure can reduce work required before execution begins.
\item But modern processors often pipeline decode aggressively, so ISA complexity does not directly equal a slow clock.
\item Extra pipeline stages can raise branch-misprediction penalties even when frequency increases.
\item Again, one design variable cannot be considered in isolation.
\end{itemize}
\end{frame}

% 24
\begin{frame}{The Iron Law Applied to ISA Design}
\begin{center}
{\Large \iron}
\end{center}
\begin{itemize}
\item \key{Instruction count (IC):} often favors richer instructions or denser encodings.
\item \key{CPI:} often benefits from regular instructions and efficient execution pipelines.
\item \key{Cycle time:} may benefit from simpler decode/control paths.
\item ISA design can improve one term while harming another.
\item Compiler quality and microarchitecture determine where the real balance lands.
\item The Iron Law is the correct framework for performance arguments about RISC and CISC.
\end{itemize}
\end{frame}

% 25
\begin{frame}{Worked Performance Comparison}
\begin{block}{Machine R}
RISC-style program: 1.30 billion instructions, CPI 1.0, 3.0 GHz
\end{block}
\begin{block}{Machine C}
CISC-style program: 0.90 billion instructions, CPI 1.4, 2.8 GHz
\end{block}
\begin{align*}
T_R &= \frac{1.30\times 10^9\times 1.0}{3.0\times 10^9}=0.433\text{ s}\\
T_C &= \frac{0.90\times 10^9\times 1.4}{2.8\times 10^9}=0.450\text{ s}
\end{align*}
\centering \key{Lower instruction count did not guarantee lower execution time.}
\end{frame}

% 26
\begin{frame}{Code Density and Memory Footprint}
\begin{itemize}
\item Code density measures how much useful program behavior fits into a given number of instruction bytes.
\item Variable-length CISC encodings can represent common operations very compactly.
\item Fixed-width RISC instructions may consume more bytes for some instruction sequences.
\item Larger code can increase instruction-cache pressure and memory bandwidth demand.
\item ARM historically addressed code density partly through Thumb and Thumb-2 encodings; AArch64 uses fixed 32-bit instructions.
\item Code density shows that regular hardware-friendly encoding has a real system-level trade-off.
\end{itemize}
\end{frame}

% 27
\begin{frame}{Instruction Fetch and Decode Complexity}
\begin{columns}[T,onlytextwidth]
\column{0.49\textwidth}
\begin{block}{Regular fixed-width ISA}
\begin{itemize}
\item Boundaries known immediately
\item Predictable field locations
\item Easier parallel decode
\item Simpler alignment logic
\end{itemize}
\end{block}
\column{0.49\textwidth}
\begin{block}{Variable-length rich ISA}
\begin{itemize}
\item Boundaries must be discovered
\item Prefixes may alter meaning
\item Fields vary by instruction
\item More front-end logic required
\end{itemize}
\end{block}
\end{columns}
\vspace{0.4em}
Modern x86 processors invest substantial hardware in the front end precisely because compatibility with variable-length x86 instructions is valuable.
\end{frame}

% 28
\begin{frame}{ISA Design and Pipelining}
\begin{itemize}
\item Pipelines work best when instruction processing can be divided into repeatable stages.
\item RISC regularity historically made five-stage pipelines comparatively straightforward to build.
\item Separate load/store operations map naturally onto a dedicated memory-access stage.
\item Fixed instruction length simplifies fetch and decode timing.
\item Complex instructions can still be pipelined, but often require decomposition or specialized paths.
\item Modern high-performance processors go far beyond classic pipelines, yet front-end regularity still matters.
\end{itemize}
\end{frame}

% 29
\begin{frame}{ISA Design and Compiler Optimization}
\begin{itemize}
\item Compilers transform high-level code into instruction sequences that exploit the ISA efficiently.
\item RISC designs expose simple operations that compilers can reorder and combine.
\item A large register set gives compilers more places to keep temporary values.
\item CISC instructions may encode useful compound operations that reduce code size or instruction count.
\item Modern compilers use detailed processor cost models rather than assuming one ISA philosophy is always superior.
\item ISA quality therefore includes how well its features can be exploited by real compiler toolchains.
\end{itemize}
\end{frame}

% 30
\begin{frame}{RISC vs. CISC: Energy and Thermal Trade-offs}
\begin{columns}[T,onlytextwidth]
\column{0.48\textwidth}
\begin{block}{RISC tendency}
\begin{itemize}
\item Regular encodings can simplify fetch, alignment, and decode.
\item Simpler architectural operations can reduce front-end work per instruction.
\item Explicit load/store behavior can make data movement more visible and predictable.
\end{itemize}
\end{block}
\column{0.48\textwidth}
\begin{block}{CISC tendency}
\begin{itemize}
\item Variable-length or richer instructions can require more complex decode machinery.
\item One architectural instruction may perform several operations or memory accesses.
\item Better code density can reduce instruction-fetch traffic and cache pressure.
\end{itemize}
\end{block}
\end{columns}
\vspace{0.15cm}
\begin{block}{Architectural implication}
ISA style influences energy, but implementation and workload determine total efficiency.
\end{block}
\end{frame}

% 31
\begin{frame}{Energy per Instruction Is Not Energy per Program}
\begin{block}{A useful first-order model}
\[
E_{\text{program}} \approx IC \times E_{\text{instruction}}
\]
\end{block}
\begin{itemize}
\item A simpler instruction may consume less energy each time it is fetched, decoded, and executed.
\item But a RISC sequence may require more architectural instructions to complete the same task.
\item A denser CISC sequence may reduce instruction-cache accesses and front-end traffic.
\item Memory behavior can dominate energy when instructions cause cache misses or DRAM accesses.
\item The relevant metric is therefore total energy for the workload, not energy for one instruction.
\item This mirrors the Iron Law: optimizing one factor in isolation can produce the wrong conclusion.
\end{itemize}
\end{frame}

% 32
\begin{frame}{Power Becomes Heat --- and Heat Requires Cooling}
\small
\begin{columns}[T,onlytextwidth]
\column{0.46\textwidth}
\begin{block}{Energy and power}
\[
P = \frac{E}{t}
\]
Higher switching activity, voltage, frequency, and active hardware structures increase power demand.
\end{block}
\column{0.50\textwidth}
\begin{block}{System consequence}
\centering
architectural complexity\\
$\downarrow$\\
transistor switching\\
$\downarrow$\\
power consumption\\
$\downarrow$\\
heat generation\\
$\downarrow$\\
cooling requirement
\end{block}
\end{columns}
\vspace{0.12cm}
\begin{itemize}
\item Cooling may range from passive heat spreaders to fans, liquid cooling, and data-center HVAC.
\item Thermal limits can force frequency reduction, so power can ultimately constrain performance.
\end{itemize}
\end{frame}

% 33
\begin{frame}{Performance per Watt: The More Useful Question}
\small
\begin{itemize}
\item Processor design increasingly optimizes useful work per unit of energy, not performance alone.
\item Simpler decode and regular execution can help efficiency, especially in constrained systems.
\item Code density, cache behavior, speculation, vector units, and memory traffic can reverse simple ISA-level expectations.
\item Process technology, voltage, frequency, and microarchitecture often matter more than the RISC/CISC label.
\item ARM's success now spans phones, laptops, and cloud servers, reflecting sustained attention to efficient implementations and system design.
\item Therefore, ``ARM is efficient because RISC is simpler'' is an incomplete architectural explanation.
\end{itemize}
\begin{block}{Architectural implication}
Compare \key{energy per workload}, \key{performance per watt}, and \key{thermal constraints}; ISA labels are only part of the evidence.
\end{block}
\end{frame}

% 31
\begin{frame}{Why Modern x86 Looks RISC-Like Inside}
\begin{itemize}
\item Modern x86 cores preserve the complex x86 ISA for software compatibility.
\item The front end decodes many x86 instructions into simpler internal micro-operations (\(\mu\)ops).
\item Those \(\mu\)ops are scheduled onto regular execution resources inside the core.
\item This separates a complex architectural interface from a more regular internal execution model.
\item Some x86 instructions decode into one \(\mu\)op; others require several or specialized handling.
\item The lesson: \key{ISA style and microarchitecture style are not the same thing.}
\end{itemize}
\end{frame}

% 32
\begin{frame}{Why Modern RISC ISAs Are Not Necessarily ``Simple''}
\begin{itemize}
\item Modern RISC ISAs include SIMD, cryptography, virtualization, atomics, and system-control instructions.
\item AArch64 has many instruction classes and sophisticated addressing features.
\item High-performance ARM cores use branch prediction, out-of-order execution, speculation, and large caches.
\item None of that violates the central RISC ideas of regular encoding and load/store organization.
\item ``Reduced'' should be understood historically as architectural regularity, not a tiny instruction count.
\item Calling a modern ARM core ``simple'' would confuse ISA philosophy with implementation sophistication.
\end{itemize}
\end{frame}

% 33
\begin{frame}{RISC Beyond Phones: Laptops}
\begin{itemize}
\item Modern RISC is no longer confined to embedded controllers or mobile phones.
\item Apple Silicon brought Arm-based processors into mainstream laptop and desktop computing.
\item Current MacBook Pro systems use Apple M5-family chips for high-performance portable computing.
\item Apple emphasizes both performance and battery life; its MacBook Pro overview lists up to 24 hours of battery life.
\item This makes laptops a useful example of the performance-per-watt argument, not just the low-power argument.
\item The architectural lesson is that RISC principles can scale from constrained devices to high-performance personal computers.
\end{itemize}
\vspace{0.15cm}
{\scriptsize Source: Apple MacBook Pro overview and technical specifications.}
\end{frame}

% 34
\begin{frame}{RISC in the Data Center: Google Axion}
\begin{itemize}
\item Cloud providers now deploy Arm-based processors for general-purpose server workloads.
\item Google Axion is a custom Arm-based CPU designed for Google Cloud data-center workloads.
\item Google identifies web and app servers, databases, analytics, media processing, and CPU-based AI as Axion target workloads.
\item Google reports production deployment for services such as Bigtable, Spanner, BigQuery, Blobstore, Pub/Sub, Google Earth Engine, and YouTube Ads.
\item These are not small embedded workloads; they are latency-sensitive and throughput-oriented cloud services.
\item The architectural question becomes performance per watt, fleet-level cost, cooling, and software ecosystem support.
\end{itemize}
\vspace{0.15cm}
{\scriptsize Source: Google Cloud Axion announcement and C4A deployment notes.}
\end{frame}

% 35
\begin{frame}{Data-Center Workloads: What Runs on Arm?}
\begin{columns}[T,onlytextwidth]
\column{0.48\textwidth}
\begin{block}{Storage and databases}
\begin{itemize}
\item Spanner
\item Bigtable
\item Cloud SQL / AlloyDB-class workloads
\end{itemize}
\end{block}
\begin{block}{Analytics and pipelines}
\begin{itemize}
\item BigQuery
\item Dataflow-style pipelines
\item batch and containerized jobs
\end{itemize}
\end{block}
\column{0.48\textwidth}
\begin{block}{Serving and platforms}
\begin{itemize}
\item YouTube Ads platform
\item web and application servers
\item containerized microservices
\end{itemize}
\end{block}
\begin{block}{AI and inference}
\begin{itemize}
\item CPU-based inference
\item preprocessing and data movement
\item services where accelerator use is not dominant
\end{itemize}
\end{block}
\end{columns}
\end{frame}

% 36
\begin{frame}{ARM AArch64 as a Modern RISC ISA}
\begin{itemize}
\item AArch64 uses fixed 32-bit instruction encodings.
\item Most arithmetic and logical operations use register operands.
\item Memory is accessed primarily through explicit load and store instructions.
\item Thirty-one general-purpose integer registers support register-oriented computation.
\item Instruction formats are structured to make common fields relatively regular.
\item These properties make AArch64 an excellent architecture for studying the connection between ISA and datapath design.
\end{itemize}
\begin{block}{Preview}
Next we move from philosophy to the actual AArch64 programmer's model and assembly syntax.
\end{block}
\end{frame}

% 34
\begin{frame}{RISC vs. CISC: What Actually Matters Today?}
\begin{itemize}
\item The binary labels are less predictive than they were in the 1980s.
\item Modern CISC processors use sophisticated internal decomposition and execution engines.
\item Modern RISC processors contain highly complex microarchitectures and rich ISA extensions.
\item Compatibility, code density, compiler ecosystem, implementation quality, and workload matter enormously.
\item The enduring RISC contribution is emphasis on regularity, load/store organization, and compiler-visible simple operations.
\item The enduring lesson is to evaluate architectural choices quantitatively rather than by slogans.
\end{itemize}
\end{frame}

% 35
\begin{frame}{Synthesis: ISA Choices Are Trade-offs}
\begin{itemize}
\item CISC tends to expose richer operations and often achieves strong code density.
\item RISC tends to expose more regular operations that are easier to decode and pipeline.
\item Neither lower instruction count nor higher clock rate guarantees better performance.
\item The Iron Law forces us to consider instruction count, CPI, and cycle time together.
\item Modern processors blur the historical boundary while preserving very different architectural interfaces.
\item AArch64 gives us a modern RISC case study for the rest of the course.
\end{itemize}
\begin{block}{Design question}
Do not ask ``Which philosophy is faster?'' Ask: \key{what trade-off does this ISA choice create, and what does the measured workload say?}
\end{block}
\end{frame}

\end{document}
